\documentclass{article}
\usepackage[T1]{fontenc}
\usepackage{lmodern}
\usepackage{graphicx}
\usepackage{amsmath,amssymb}
\usepackage[a4paper,margin=2cm]{geometry}
\usepackage[absolute,overlay]{textpos}
\setlength{\TPHorizModule}{1mm}
\setlength{\TPVertModule}{1mm}
\usepackage[indonesian]{babel}
\usepackage{hyperref}
\setlength{\emergencystretch}{2em}
% unit: Halaman 57

\begin{document}

% === Halaman 57 (llm) ===

\textbf{Contoh 2}: Misalkan $a = 01010111 = x^6 + x^4 + x^2 + x + 1$ dan $b = 10000011 = x^7 + x + 1$
$a$ dan $b \in GF(2^8)$ \quad (dalam notasi Hex, $a = '57'$ dan $b = '83'$)

(i) $\begin{align*} a + b &= '57' + '83' = (x^6 + x^4 + x^2 + x + 1) + (x^7 + x + 1) \\ &= x^7 + x^6 + x^4 + x^2 + 2x + 2 \pmod{2} \\ &= x^7 + x^6 + x^4 + x^2 + 0x + 0 \\ &= x^7 + x^6 + x^4 + x^2 = 11010100 = 'D4' \end{align*}$ \vspace{5mm} $\text{\textcolor{red}{Bagi tiap koefisien dengan 2, lalu ambil sisanya}}$

Dalam representasi biner:

$\begin{array}{rl} 01010111 \\ 10000011 + \\ \hline 11010100 \end{array} \rightarrow \text{sama dengan hasil operasi XOR}$

$\therefore a + b = 11010100 = a \text{ XOR } b$

\end{document}
